\documentclass{beamer}
\usetheme{UASLP1}
\usepackage[utf8]{inputenc}
\usepackage[T1]{fontenc}
\usepackage[spanish]{babel}
\usepackage{helvet}
\usepackage{bookmark}
\usepackage{amsmath}
\usepackage{amssymb}
\usepackage{bm}
\usepackage[useregional]{datetime2}
\usepackage{booktabs}
\usepackage{caption}
\usepackage{mathtools}

\renewcommand{\today}{\ifcase\month\or
  Enero\or Febrero\or Marzo\or Abril\or Mayo\or Junio\or
  Julio\or Agosto\or Septiembre\or Octubre\or Noviembre\or Diciembre\fi\space \number\year}

%--- INFORMACIÓN DE LA PORTADA ---
\title{Marco Teórico: Identificación de Modos Electromecánicos}
\subtitle{\textit{Algoritmo Multi-TK }}
\author{Ing. Glader David Hernández Reyes}
\date{Septiembre 2025}
\institute{Director: Dr. José de Jesús Nuño Ayón}

\begin{document}

%--- PORTADA ---
\begin{frame}[plain,t]
\titlepage
\end{frame}

%--- ÍNDICE ---
\begin{frame}
\frametitle{Contenido}
\tableofcontents
\end{frame}

% ===========================================================
\section{Dinámica de Oscilaciones Electromecánicas}
% ===========================================================

\begin{frame}
\frametitle{Dinámica de Oscilaciones Electromecánicas}
\framesubtitle{Ecuaciones de Swing del Generador Síncrono}

La dinámica de cada generador $i$ en un sistema de $n_g$ máquinas ($i=1,\ldots,n_g$) se describe mediante las \textbf{ecuaciones de swing}:

\begin{block}{Ecuaciones de Swing }
\begin{align}
    \dot{\delta}_i  &= \omega_i - \omega_s \label{eq:swing1} \\[4pt]
    \dot{\omega}_i  &= \frac{\omega_s}{2H_i}
                       \Bigl(P_{m,i} - P_{e,i}(\boldsymbol{\delta}) - D_i(\omega_i - \omega_s)\Bigr)
                       \label{eq:swing2}
\end{align}
\end{block}

Para generadores síncronos modeloados de forma clásica (voltaje constante a partir de una reactancia transitoria). 
\end{frame}

% -----------------------------------------------------------
\begin{frame}
Parámetros: 
\vspace{0.2em}
\begin{itemize}
    \item $\delta_i$: ángulo del rotor (rad);
    \item $\omega_i$: velocidad angular del rotor (rad/s);
    \item $H_i$: constante de inercia (s); 
    \item $D_i$: constante de amortiguamiento (s);
    \item $P_{m,i}$: potencia mecánica; 
    \item $P_{e,i}(\boldsymbol{\delta})$: potencia eléctrica — función no lineal de los ángulos de rotor $\boldsymbol{\delta} = [\delta_1, \ldots, \delta_{n_g}]$;
    \item $\omega_s = 2\pi f_s$: velocidad síncrona de referencia
\end{itemize}
\end{frame}

% -----------------------------------------------------------


% -----------------------------------------------------------
\begin{frame}
 \frametitle{Dinámica de Oscilaciones Electromecánicas}
 \framesubtitle{Linealización alrededor del Punto de Operación}

 Para analizar la estabilidad de pequeña señal, se linealiza el sistema alrededor del punto de equilibrio
 $(\boldsymbol{\delta}^0, \boldsymbol{\omega}^0)$, en régimen permanente se cumple ($\dot{\delta}_i = 0$, $\dot{\omega}_i = 0$): 
 \begin{equation}
    P_{m,i} = P_{e,i}(\boldsymbol{\delta}^0), \qquad \omega_i^0 = \omega_s  
 \end{equation}

Se definen las variables de estado de perturbación (desviaciones):
 \begin{equation}
    \Delta\boldsymbol{\delta} = \boldsymbol{\delta} - \boldsymbol{\delta}^0, \qquad
    \Delta\boldsymbol{\omega} = \boldsymbol{\omega} - \boldsymbol{\omega}^0
 \end{equation}

\end{frame}
% -----------------------------------------------------------
\begin{frame}
    Linealizando la primera ecuacion de osilación:
    \begin{equation}
        \Delta\dot{\boldsymbol{\delta}} = \Delta\boldsymbol{\omega}
    \end{equation}

 Expandiendo la serie de Taylor de primero orden de $P_{e,i}(\boldsymbol{\delta})$ alrededor de $\boldsymbol{\delta}^0$ y usando la condicion de equilibrio se linealiza la segunda ecuacion de osilación:
 \begin{equation}
    \Delta\dot{\boldsymbol{\omega}}_i = \frac{\omega_s}{2H_i}\Bigl(-\sum_{j=1}^{n_g} \frac{\partial P_{e,i}}{\partial \delta_j}\big|_0 \Delta\delta_j - D_i \Delta\omega_i\Bigr)
  \end{equation}
    
\end{frame}
% -----------------------------------------------------------


% -----------------------------------------------------------
\begin{frame}
Definimos los vectores de desviaciones: 
\begin{equation}
    \Delta\boldsymbol{\delta} = [\Delta\delta_1, \ldots, \Delta\delta_{n_g}]^\top, \qquad
    \Delta\boldsymbol{\omega} = [\Delta\omega_1, \ldots, \Delta\omega_{n_g}]^\top
\end{equation}
 las matrices diagonales:
 \begin{equation}
    \mathbf{M} = \mathrm{diag}\left(\frac{2H_i}{\omega_s}\right), \qquad
    \mathbf{K_D} = \mathrm{diag}(D_i)
 \end{equation}
 La matriz de sincronización con la red (o matriz de coeficientes de par sincronizante) representa el cambio en las potencias eléctricas para un cambio en el ángulo en el punto de equilibrio: 
\begin{equation}
    \mathbf{K}_s = \left[\frac{\partial P_{e,i}}{\partial \delta_j}\big|_{\delta = \delta^0}\right]_{i,j=1}^{n_g}
\end{equation}

\end{frame}
% -----------------------------------------------------------

% -----------------------------------------------------------
\begin{frame}
Con estas definiciones y las ecuaciones linealizadas podemos definir la ecuacion en forma matricial: 

       \begin{equation}
    \begin{bmatrix} \Delta\dot{\boldsymbol{\delta}} \\ \Delta\dot{\boldsymbol{\omega}} \end{bmatrix}
    =
    \underbrace{
    \begin{bmatrix}
        \mathbf{0}  & \mathbf{I} \\
        -\mathbf{M}^{-1}\mathbf{K}_s & -\mathbf{M}^{-1}\mathbf{K_D}
    \end{bmatrix}
    }_{\mathbf{A} \in \mathbb{R}^{2n_g \times 2n_g}}
    \begin{bmatrix} \Delta\boldsymbol{\delta} \\ \Delta\boldsymbol{\omega} \end{bmatrix}
    \label{eq:linear_sys}
\end{equation}
Para un sistema autónomo ($u=0$) se tiene una solución de la forma: 
\begin{equation}
    \begin{bmatrix} \Delta\boldsymbol{\delta}(t) \\ \Delta\boldsymbol{\omega}(t) \end{bmatrix}
    = e^{\mathbf{A}t}\,
    \begin{bmatrix} \Delta\boldsymbol{\delta}(0) \\ \Delta\boldsymbol{\omega}(0) \end{bmatrix}
    \label{eq:solution}
\end{equation}


\end{frame}
% -----------------------------------------------------------
\begin{frame}

Suponiendo que la matriz $\mathbf{A}$ es diagonalizable (lo que ocurre en la práctica salvo casos muy degenerados), podemos escribir
\[
\mathbf{A} = \mathbf{V} \mathbf{\lambda} \mathbf{W}^\top, \quad \mathbf{W}^\top = \mathbf{V}^{-1},
\]
donde
\begin{itemize}
    \item $\mathbf{\lambda} = \text{diag}(\lambda_1, \lambda_2, \dots, \lambda_{2n_g})$ contiene los \textbf{eigenvalores}.
    \item $\mathbf{V} = [ \mathbf{v}_1 \ \mathbf{v}_2 \ \dots \ \mathbf{v}_{2n_g} ]$ es la matriz de \textbf{eigenvectores derechos} (cada $\mathbf{v}_k$ es un vector columna de tamaño $2n_g$).
    \item $\mathbf{W} = [ \mathbf{w}_1 \ \mathbf{w}_2 \ \dots \ \mathbf{w}_{2n_g} ]$ contiene los \textbf{eigenvectores izquierdos} (vectores columna), de modo que $\mathbf{w}_k^\top$ es una fila de $\mathbf{V}^{-1}$.
\end{itemize}

\end{frame}


\begin{frame}

La respuesta libre a partir de una condición inicial $\mathbf{x}(0)$ (la perturbación) se expresa como suma de modos:
\begin{equation}
\mathbf{x}(t) = e^{\mathbf{A}t}\mathbf{x}(0) = \sum_{k=1}^{2n_g} \mathbf{v}_k \left( \mathbf{w}_k^\top \mathbf{x}(0) \right) e^{\lambda_k t}.
\end{equation}

Esta ecuación es fundamental: muestra que el estado en cualquier instante es la superposición de los modos, cada uno con una ``intensidad'' $\alpha_k = \mathbf{w}_k^\top \mathbf{x}(0)$ (un escalar complejo si el modo es complejo; para modos reales sería real).


\end{frame}

\begin{frame}
La ecuación de salida:
\[
y(t) = \mathbf{c}^\top \mathbf{x}(t) = \sum_{k=1}^{2n_g} \left( \mathbf{c}^\top \mathbf{v}_k \right) \left( \mathbf{w}_k^\top \mathbf{x}(0) \right) e^{\lambda_k t}.
\]

Definimos el \textbf{residuo complejo} asociado al modo $k$ como
\begin{equation}
\boxed{R_k = \left( \mathbf{c}^\top \mathbf{v}_k \right) \left( \mathbf{w}_k^\top \mathbf{x}(0) \right)}.
\end{equation}

Con esto,
\[
y(t) = \sum_{k=1}^{2n_g} R_k e^{\lambda_k t}.
\]


\end{frame}

\begin{frame}
    Dado que los eigenvalores aparecen en pares complejos conjugados ($\lambda_{k+1} = \lambda_k^*$) y los correspondientes $R_k$ también son conjugados, podemos agruparlos para escribir la típica sinusoide amortiguada:
\[
y(t) = \sum_{\text{modos reales}} R_k e^{\sigma_k t} + \sum_{\text{pares conjugados}} 2 |R_k| e^{\sigma_k t} \cos(\omega_k t + \angle R_k).
\]

En el contexto de identificación, se suele denotar $R_{ik}$ para el canal $i$ (aquí $R_k$ para un solo canal).



\end{frame}

\begin{frame}

Cada par conjugado $\lambda_k = \sigma_k \pm j\omega_k$ genera en el tiempo
una componente oscilatoria:
\[
y_k(t) = 2|R_k| e^{\sigma_k t} \cos(\omega_k t + \phi_k).
\]
\begin{itemize}
    \item $\omega_k$: frecuencia angular de la oscilación [rad/s].
    \item Frecuencia en \textbf{Hz}: $f_k = \dfrac{\omega_k}{2\pi}$.
    \item El decaimiento está gobernado por $\sigma_k$; para cuantificarlo
          sin unidades se usa la \textbf{razón de amortiguamiento} $\zeta_k$,
          que se define igualando el polinomio característico del modo con
          el de un oscilador de segundo orden:
          \[
          s^2 + 2\zeta_k\omega_{n,k} s + \omega_{n,k}^2 = 0,
          \]
          donde $\omega_{n,k} = \sqrt{\sigma_k^2+\omega_k^2}$.
          
\end{itemize}
\end{frame}

% -----------------------------------------------------------
\begin{frame}
\frametitle{Dinámica de Oscilaciones Electromecánicas}
\framesubtitle{Eigenvalores y Modos del Sistema}

%a respuesta del sistema linealizado está gobernada por los \textbf{eigenvalores} de $\mathbf{A}$:

\begin{equation}
    \lambda_k = \sigma_k \pm j\omega_k, \qquad k = 1,\ldots,n_g
    \label{eq:eigenvalue}
\end{equation}

Los parámetros modales asociados a cada par complejo conjugado son:

\begin{columns}
\column{0.48\textwidth}
\begin{block}{Frecuencia de Oscilación}
\begin{equation}
    f_k = \frac{\omega_k}{2\pi} \quad [\text{Hz}]
    \label{eq:freq}
\end{equation}
\end{block}
\begin{block}{Razón de Amortiguamiento}
\begin{equation}
    \zeta_k = \frac{-\sigma_k}{\sqrt{\sigma_k^2 + \omega_k^2}}
    \label{eq:damp}
\end{equation}
\end{block}

\column{0.48\textwidth}
\begin{alertblock}{Condición de Estabilidad}
    El sistema es estable si y solo si:
    \begin{equation}
        \sigma_k < 0 \quad \forall\, k
        \label{eq:stability}
    \end{equation}
    Convencionalmente se requiere $\zeta_k \geq 5\%$.
\end{alertblock}
\end{columns}

\end{frame}

% -----------------------------------------------------------
\begin{frame}
\frametitle{Dinámica de Oscilaciones Electromecánicas}
\framesubtitle{Respuesta Temporal: Modelo de Exponenciales Amortiguadas}

La respuesta medida (velocidad angular, ángulo, potencia) de cada señal $i$ tras una perturbación puede expresarse como:

\begin{equation}
    y_i(t) = \sum_{k=1}^{n} A_{ik}\, e^{\sigma_k t}\cos(\omega_k t + \phi_{ik})
           = \mathrm{Re}\!\left\{\sum_{k=1}^{n} R_{ik}\, e^{\lambda_k t}\right\}
    \label{eq:response}
\end{equation}

donde $R_{ik} = A_{ik}e^{j\phi_{ik}}$ es el \textbf{residuo complejo} del modo $k$ en la señal $i$.

\vspace{0.4em}
\begin{itemize}
    \item $n$: número de modos excitados por la perturbación
    \item $A_{ik}$: amplitud modal (magnitud del residuo)
    \item $\phi_{ik}$: fase modal — define la \textbf{forma modal} $\bm{\phi}_k$
\end{itemize}



\end{frame}


\begin{frame}
    \begin{exampleblock}{Objetivo de la estimación modal}
    Recuperar $\{\lambda_k, R_{ik}\}_{k=1}^n$ directamente desde las mediciones $y_i(t)$,\textbf{ sin requerir el modelo} $\mathbf{A}$.
\end{exampleblock}
\end{frame}



% ===========================================================
\section{Formulación Discreta y Predicción Lineal}
% ===========================================================

\begin{frame}
\frametitle{Formulación Discreta y Predicción Lineal}
\framesubtitle{Del Tiempo Continuo al Dominio Muestreado}

Muestreando cada señal $x_i(t)$ con periodo $\Delta t$, definimos
$x_i[n] = x_i(n\Delta t)$, para $n = 0,1,\dots,N-1$. 
Bajo el modelo de suma de exponenciales amortiguadas,
\begin{equation}
    x_i[n] = \sum_{k=1}^{M} a_{i,k}\, z_k^{\,n}, \qquad
    z_k = e^{\lambda_k \Delta t} = e^{(-\delta_k + j\omega_k)\Delta t}.
    \label{eq:discrete}
\end{equation}

\end{frame}

\begin{frame}

Los polos en tiempo discreto $z_k$ contienen toda la información modal:

\begin{block}{Mapeo Continuo $\leftrightarrow$ Discreto}
\begin{align}
    z_k &= e^{(-\delta_k + j\omega_k)\Delta t}, \\[4pt]
    \lambda_k = -\delta_k + j\omega_k &= \frac{\ln z_k}{\Delta t}.
    \label{eq:mapping}
\end{align}
\end{block}

\vspace{0.3em}
Esta relación permite recuperar el amortiguamiento $\delta_k$ y la frecuencia
$\omega_k$ a partir de las raíces $z_k$ del polinomio de predicción.

\end{frame}

\begin{frame}
\frametitle{Formulación Discreta y Predicción Lineal}
\framesubtitle{Del Tiempo Continuo al Dominio Muestreado}

Muestreando $x_i(t)$ con periodo $\Delta t$, la muestra $x_i[n] = x_i(n\Delta t)$ satisface:

\begin{equation}
    x_i[n] = \sum_{k=1}^{M} a_{i,k}\, z_k^n, \qquad z_k = e^{\lambda_k \Delta t}
    \label{eq:discrete}
\end{equation}

\vspace{0.3em}
La secuencia \eqref{eq:discrete} satisface una \textbf{recurrencia lineal de orden $L$}:

\begin{equation}
    x_i[n] + \sum_{\ell=1}^{L} b_\ell\, x_i[n-\ell] = 0, \qquad n \geq L
    \label{eq:ar}
\end{equation}

donde los coeficientes $\{b_\ell\}$ son las componentes del \textbf{polinomio de predicción} cuyas raíces son exactamente $\{z_k\}$.

\end{frame}

% -----------------------------------------------------------
\begin{frame}
\frametitle{Formulación Discreta y Predicción Lineal}
\framesubtitle{Extracción de Parámetros Modales}

Una vez obtenidos los coeficientes $\{b_\ell\}_{\ell=1}^{L}$, los polos discretos se obtienen resolviendo el polinomio característico de orden $L$:

\begin{equation}
    P(z) = z^{L} + b_1 z^{L-1} + \cdots + b_L = 0
    \quad \Rightarrow \quad \{z_k\}_{k=1}^{L}
    \label{eq:poly}
\end{equation}

De estas $L$ raíces, solo las $M$ asociadas a la señal corresponden a modos electromecánicos.

La relación con los polos continuos $s_k = -\delta_k + j\omega_k$ está dada por $z_k = e^{s_k \Delta t}$, lo que permite extraer los parámetros:

\begin{align}
    \delta_k &= -\frac{1}{2\Delta t}\,\ln\!\bigl([\Re(z_k)]^2 + [\Im(z_k)]^2\bigr) \label{eq:delta} \\
    \omega_k &= \frac{1}{\Delta t}\,\arg(z_k) \label{eq:omega}
\end{align}

\end{frame}


\begin{frame}
 Equivalentemente, si se introduce $\lambda_k = \frac{\ln z_k}{\Delta t} = -\delta_k + j\omega_k$, entonces
 $\delta_k = -\Re(\lambda_k)$ y $\omega_k = \Im(\lambda_k)$.
\end{frame}

\begin{frame}
    \frametitle{Formulación Discreta y Predicción Lineal}
    \framesubtitle{Condición de estabilidad y parámetros finales}

    \begin{alertblock}{Condición de estabilidad en tiempo discreto}
        El sistema es estable si y solo si todos los polos $z_k$ están \textbf{dentro} del círculo unitario:
        \begin{equation}
            |z_k| < 1 \quad \forall\, k
            \label{eq:stability_z}
        \end{equation}
        lo que equivale a $\delta_k > 0$ para cada modo real.
    \end{alertblock}

    \vspace{0.3em}
    Una vez calculados $\delta_k$ y $\omega_k$, los parámetros modales característicos son:

    \begin{equation}
    f_k = \dfrac{\omega_k}{2\pi}\quad [Hz].
    \label{eq:freq_final}
    \end{equation}
    \begin{equation}
    \zeta_k = \dfrac{\delta_k}{\sqrt{\delta_k^2 + \omega_k^2}} \quad ( \zeta_k\% = \zeta_k \times 100).
    \label{eq:damp_final}
    \end{equation}

\end{frame}

% ===========================================================
\section{Algoritmo Tufts-Kumaresan Multivariable (Multi-TK)}
% ===========================================================

\begin{frame}
\frametitle{Algoritmo Multi-TK}
\framesubtitle{Modelo de señal y predicción lineal hacia atrás}

Para $P$ señales medidas $\{x_i[n]\}_{i=1}^P$, cada una de $N$ muestras ($n=0,\dots,N-1$), el modelo de señal es
\[
x_i[n] = \sum_{k=1}^{M} a_{i,k}\, z_k^n + w_i[n], \qquad z_k = e^{(-\delta_k + j\omega_k)\Delta t}.
\]

Se elige un orden de predictor $L$ tal que $M \le L \le N-M$. Para cada canal $i$ se plantea el problema de \textbf{predicción lineal hacia atrás}:
\begin{equation}
\mathbf{A}_i\,\mathbf{b} = -\mathbf{h}_i,
\label{eq:backward_pred}
\end{equation}
donde $\mathbf{b}=[b(1),b(2),\dots,b(L)]^\top$ es el vector de coeficientes del polinomio característico.
\end{frame}


\begin{frame}
\frametitle{Algoritmo Multi-TK}

La matriz de Hankel conjugada para la $i$-ésima señal se construye como:
\begin{equation}
\mathbf{A}_i =
\begin{bmatrix}
x_i^*(1)   & x_i^*(2)   & \cdots & x_i^*(L) \\
x_i^*(2)   & x_i^*(3)   & \cdots & x_i^*(L+1) \\
\vdots     & \vdots     & \ddots & \vdots   \\
x_i^*(N-L) & x_i^*(N-L+1) & \cdots & x_i^*(N-1)
\end{bmatrix}
\in \mathbb{C}^{(N-L)\times L},
\label{eq:Ai}
\end{equation}
y el vector objetivo asociado es
\begin{equation}
\mathbf{h}_i =
\begin{bmatrix}
x_i^*(0) \\ x_i^*(1) \\ \vdots \\ x_i^*(N-L-1)
\end{bmatrix}
\in \mathbb{C}^{N-L}.
\label{eq:hi}
\end{equation}
\end{frame}

\begin{frame}
\frametitle{Algoritmo Multi-TK}
\framesubtitle{Sistema global multicanal}

Se apilan verticalmente las matrices y los vectores de todos los canales:
\begin{equation}
\mathbf{A}_{\text{multi}} =
\begin{bmatrix} \mathbf{A}_1 \\ \mathbf{A}_2 \\ \vdots \\ \mathbf{A}_P \end{bmatrix}
\in \mathbb{C}^{P(N-L)\times L},
\qquad
\mathbf{h}_{\text{multi}} =
\begin{bmatrix} \mathbf{h}_1 \\ \mathbf{h}_2 \\ \vdots \\ \mathbf{h}_P \end{bmatrix}
\in \mathbb{C}^{P(N-L)}.
\label{eq:multi_system}
\end{equation}

El sistema completo de predicción lineal es:
\begin{equation}
\boxed{\mathbf{A}_{\text{multi}}\,\mathbf{b} = -\mathbf{h}_{\text{multi}}}.
\label{eq:multi_equation}
\end{equation}
\end{frame}
% -----------------------------------------------------------
\begin{frame}
\frametitle{Algoritmo Multi-TK}
\framesubtitle{Descomposición SVD y selección del orden del modelo}

Se aplica la \textbf{descomposición en valores singulares} a $\mathbf{A}_{\text{multi}}$:
\begin{equation}
\mathbf{A}_{\text{multi}} = \mathbf{U}\,\boldsymbol{\Sigma}\,\mathbf{V}^H,
\label{eq:svd}
\end{equation}
con $\boldsymbol{\Sigma} = \operatorname{diag}(\sigma_1,\sigma_2,\dots,\sigma_r)$, $\sigma_1 \ge \sigma_2 \ge \dots \ge \sigma_r > 0$.

El número de modos $M$ se selecciona mediante el criterio de \textbf{energía acumulada}:
\begin{equation}
E_{s,M} = \frac{\displaystyle\sum_{k=1}^{M} \sigma_k^2}{\displaystyle\sum_{k=1}^{r} \sigma_k^2} \ge E_0,
\label{eq:energy}
\end{equation}
donde típicamente $E_0 = 97\%$ (o $99.9\%$ según el nivel de ruido).
\end{frame}

% -----------------------------------------------------------
\begin{frame}
\frametitle{Algoritmo Multi-TK}
\framesubtitle{Truncamiento de la SVD y solución del sistema}

Se retiene el subespacio de señal de rango $M$:
\begin{equation}
\mathbf{A}_{\text{multi}} \approx \mathbf{U}_s \boldsymbol{\Sigma}_s \mathbf{V}_s^H,
\quad
\mathbf{U}_s \in \mathbb{C}^{P(N-L)\times M},\;
\boldsymbol{\Sigma}_s \in \mathbb{R}^{M\times M},\;
\mathbf{V}_s \in \mathbb{C}^{L\times M}.
\label{eq:truncsvd}
\end{equation}

La solución regularizada del sistema $\mathbf{A}_{\text{multi}}\,\mathbf{b} = -\mathbf{h}_{\text{multi}}$ es:
\begin{equation}
\boxed{\mathbf{b} = -\mathbf{V}_s\,\boldsymbol{\Sigma}_s^{-1}\,\mathbf{U}_s^H\,\mathbf{h}_{\text{multi}}}.
\label{eq:solution}
\end{equation}

\begin{itemize}
\item El truncamiento elimina componentes de ruido y reduce el sesgo en la estimación de $\mathbf{b}$.
\item Orden del predictor: se recomienda $L \approx \tfrac{3N}{4}$.
\end{itemize}
\end{frame}



% ----------------------------------------------------------

\begin{frame}
\frametitle{Algoritmo Multi‑TK}
\framesubtitle{Solución del Sistema de Predicción Lineal}

Dado el sistema global $\mathbf{A}_{\text{multi}}\mathbf{b} = -\mathbf{h}_{\text{multi}}$,
el vector de coeficientes de predicción $\mathbf{b}\in\mathbb{C}^{L}$ se obtiene minimizando:
\begin{equation}
    \min_{\mathbf{b}}\; \|\mathbf{A}_{\text{multi}}\mathbf{b} + \mathbf{h}_{\text{multi}}\|_2^2
    \label{eq:ls}
\end{equation}

Tras la SVD truncada a $M$ componentes ($\mathbf{A}_{\text{multi}} \approx \mathbf{U}_s \boldsymbol{\Sigma}_s \mathbf{V}_s^H$):
\begin{align}
    \mathbf{w} &= \boldsymbol{\Sigma}_s^{-1}\,\mathbf{U}_s^H\,\mathbf{h}_{\text{multi}}
                 \label{eq:weights}\\[4pt]
    \mathbf{b} &= -\mathbf{V}_s\,\mathbf{w}
                 \label{eq:bcoeffs}
\end{align}

\vspace{0.3em}
La solución retiene solo los $M$ valores singulares dominantes, lo que reduce la influencia del ruido.
\end{frame}
\begin{frame}

\frametitle{Algoritmo Multi‑TK}

Los coeficientes $\mathbf{b}=[b_1,\dots,b_L]^\top$ definen el polinomio de  predicción:
\begin{equation}
    P(z) = z^L + b_1 z^{L-1} + b_2 z^{L-2} + \cdots + b_L = 0
    \label{eq:pred_poly}
\end{equation}

\begin{exampleblock}{Raíces y modos}
    Las $L$ raíces $\{z_k\}$ de $P(z)$ son los \textbf{polos discretos} del sistema.
    De ellas, solo $M$ corresponden a la dinámica real; el resto son polos
    espurios que se filtran posteriormente mediante el criterio de energía.
\end{exampleblock}
\end{frame}


\begin{frame}
\frametitle{Algoritmo Multi‑TK}
\framesubtitle{Órdenes del Predictor y del Modelo}

En el método Tufts‑Kumaresan se distingue entre:
\begin{itemize}
    \item \textbf{Orden del predictor} $L$: número de columnas de $\mathbf{A}_i$ y
          grado del polinomio $P(z)$. Se elige \emph{a priori} tal que
          $M \le L \le N-M$.
    \item \textbf{Número de modos} $M$: cantidad de componentes significativas
          (valores singulares conservados). Se determina \emph{automáticamente}
          mediante el criterio de energía acumulada (Ec.~\ref{eq:energy}).
\end{itemize}

\vspace{0.4em}
Elegir $L > M$ introduce grados de libertad extra que permiten modelar el
ruido sin afectar las estimaciones de los verdaderos modos; las raíces
espurias se descartan en la etapa de filtrado.
\end{frame}

\begin{frame}
\frametitle{Algoritmo Multi‑TK}
\framesubtitle{Resumen del Flujo Computacional}

\begin{enumerate}
    \item \textbf{Pre‑procesamiento:} Filtrado pasa‑banda $[f_{\min}, f_{\max}]$ y diezmado.
    \item \textbf{Apilamiento:} Construcción de $\mathbf{A}_{\text{multi}}$ y $\mathbf{h}_{\text{multi}}$
          desde las $P$ señales (Ec.~\ref{eq:multi_system}).
    \item \textbf{SVD:} $\mathbf{A}_{\text{multi}} = \mathbf{U}\boldsymbol{\Sigma}\mathbf{V}^H$.
    \item \textbf{Selección de orden:} $M$ mediante energía acumulada $\ge E_0$ (Ec.~\ref{eq:energy}).
    \item \textbf{Coeficientes:} Solución LS regularizada $\mathbf{b}$ (Ec.~\ref{eq:weights}--\ref{eq:bcoeffs}).
    \item \textbf{Polos discretos:} Raíces de $P(z) = z^L + b_1 z^{L-1} + \cdots + b_L$ (Ec.~\ref{eq:pred_poly}).
    \item \textbf{Parámetros modales:} Mapeo $z_k \to \lambda_k$ y cálculo de $\delta_k, \omega_k, f_k, \zeta_k$
          (Ec.~\ref{eq:delta}--\ref{eq:damp_final}.)
\end{enumerate}

%\vspace{0.4em}
%\begin{alertblock}{Filtrado de raíces}
 %   Solo se conservan raíces con $|z_k^{-1}| < 1$ (polos dentro del círculo unitario),
  %  lo que equivale a $\delta_k > 0$ (estabilidad en tiempo continuo).
%\end{alertblock}

\end{frame}

% -----------------------------------------------------------


% ===========================================================
\section{Patrones Dinámicos: Formas Modales y Factores de Participación}
% ===========================================================
\begin{frame}
\frametitle{Patrones Dinámicos}
\framesubtitle{Estimación de amplitudes complejas (sistema de Vandermonde)}

Conocidos los polos  $\{z_k\}_{k=1}^{M}$, para cada señal $i$ ($i=1,\dots,P$)
se resuelve el sistema lineal que relaciona las muestras con las amplitudes complejas
$a_{i,k}$:

\begin{equation}
    \underbrace{
    \begin{bmatrix}
        z_1^0 & z_2^0 & \cdots & z_M^0 \\
        z_1^1 & z_2^1 & \cdots & z_M^1 \\
        \vdots &      & \ddots & \vdots \\
        z_1^{N-1} & z_2^{N-1} & \cdots & z_M^{N-1}
    \end{bmatrix}
    }_{\boldsymbol{\Psi}\,\in\,\mathbb{C}^{N \times M}}
    \begin{bmatrix} a_{i,1} \\ a_{i,2} \\ \vdots \\ a_{i,M} \end{bmatrix}
    = \mathbf{x}_i
    \label{eq:vander_art}
\end{equation}
\end{frame}

\begin{frame}
    
La solución de mínimos cuadrados (pseudoinversa de $\boldsymbol{\Psi}$) es:

\begin{equation}
    \mathbf{a}_i = \boldsymbol{\Psi}^+\,\mathbf{x}_i
    = \bigl(\boldsymbol{\Psi}^H\boldsymbol{\Psi}\bigr)^{-1}\boldsymbol{\Psi}^H\mathbf{x}_i
    \label{eq:residues_art}
\end{equation}

De cada elemento $a_{i,k}$ se obtienen la amplitud y la fase:
\[
|a_{i,k}|,\qquad \phi_{i,k} = \arg(a_{i,k}).
\]
\end{frame}
\begin{frame}
\frametitle{Patrones Dinámicos}
\framesubtitle{Forma Modal (\textit{mode shape})}

La \textbf{forma modal} del modo $k$ describe la distribución espacial de la oscilación
a través de las $P$ señales. Se construye a partir de las amplitudes complejas:

\begin{equation}
    \mathbf{s}_k =
    \begin{bmatrix}
        a_{1,k} \\
        a_{2,k} \\
        \vdots \\
        a_{P,k}
    \end{bmatrix}
    =
    \begin{bmatrix}
        |a_{1,k}|\,e^{j\phi_{1,k}} \\
        |a_{2,k}|\,e^{j\phi_{2,k}} \\
        \vdots \\
        |a_{P,k}|\,e^{j\phi_{P,k}}
    \end{bmatrix}
    \label{eq:mode_shape_art}
\end{equation}

\begin{itemize}
    \item $|\mathbf{s}_k|$ (magnitud): indica qué tan observable es el modo $k$ en cada señal.
    \item $\angle\mathbf{s}_k$ (fase): revela las relaciones de fase entre señales;
          ángulos opuestos ($\sim 180^\circ$) indican que los generadores oscilan \textbf{en contrafase}.
\end{itemize}

\end{frame}


% -----------------------------------------------------------



\begin{frame}
\frametitle{Patrones Dinámicos}
\framesubtitle{Factor de Participación – definición basada en mediciones}

En el contexto de identificación a partir de señales, el \textbf{factor de participación}
se define a partir de la magnitud de las amplitudes complejas,
ya que éstas cuantifican la observabilidad relativa del modo en cada medición.

Para el modo $k$, la participación normalizada de la señal $i$ es:

\begin{equation}
    \boxed{\hat{p}_{ik} = \frac{|a_{i,k}|}{\displaystyle\sum_{j=1}^{P} |a_{j,k}|}}
    \label{eq:pf_mediciones}
\end{equation}

\begin{itemize}
    \item $\hat{p}_{ik} \in [0,1]$; un valor cercano a $1$ indica que el modo $k$
          es más observable en la señal $i$.
    \item Esta métrica permite localizar qué generadores o áreas dominan la oscilación.

\end{itemize}

\end{frame}


% -----------------------------------------------------------
\begin{frame}
\frametitle{Patrones Dinámicos}
\framesubtitle{Factor de Participación – Aplicaciones}

\begin{columns}
\column{0.48\textwidth}
\begin{block}{Aplicaciones}
\begin{itemize}
    \item Localización de la fuente de la oscilación.
    \item Selección óptima de la ubicación del PSS.
    \item Identificación de generadores críticos.
    \item Evaluación de observabilidad de modos en distintos puntos de la red.
\end{itemize}
\end{block}

\column{0.48\textwidth}
\begin{exampleblock}{Criterio de dominancia}
    El generador $i^*$ domina el modo $k$ si:
    \begin{equation}
        i^* = \arg\max_i \hat{p}_{ik}
        \label{eq:dominancia_art}
    \end{equation}
    \vspace{0.2em}
    Un valor alto de $\hat{p}_{ik}$ indica que la señal $i$ es la más indicada
    para monitorear o amortiguar el modo $k$.
\end{exampleblock}
\end{columns}

\end{frame}

% -----------------------------------------------------------
\begin{frame}
\frametitle{Patrones Dinámicos}
\framesubtitle{Grupos Coherentes por Coseno Director}

Dos señales $i$ y $j$ presentan \textbf{coherencia} respecto al modo $k$
si sus amplitudes complejas son aproximadamente paralelas en el plano complejo.
El grado de coherencia se mide con el \textbf{coseno director}:

\begin{equation}
    r_{ij} =
    \frac{\displaystyle\sum_{k=1}^{M}
          \bigl(a_{i,k} - \bar{a}_i\bigr)\bigl(a_{j,k} - \bar{a}_j\bigr)^*}
         {\sqrt{\displaystyle\sum_{k=1}^{M}|a_{i,k}-\bar{a}_i|^2
                \displaystyle\sum_{k=1}^{M}|a_{j,k}-\bar{a}_j|^2}}
    \label{eq:direction_cosine_art}
\end{equation}

donde $\bar{a}_i = \frac{1}{M}\sum_{k=1}^{M} a_{i,k}$.

\end{frame}


\begin{frame}
    \begin{itemize}
    \item $r_{ij} \approx +1$: señales completamente coherentes (oscilan en fase).
    \item $r_{ij} \approx -1$: señales en contrafase (grupos opuestos).
    \item $r_{ij} \approx 0$: comportamiento independiente respecto a los modos observados.
\end{itemize}
\end{frame}


% ===========================================================
\section{Resumen del Marco Teórico}
% ===========================================================

\begin{frame}
\frametitle{Resumen del Marco Teórico}
\framesubtitle{Integración de las Etapas}

\begin{enumerate}
    \item \textbf{Modelo físico:} Ecuaciones de swing (Ec.~\ref{eq:swing1}--\ref{eq:swing2})
          $\;\to\;$ respuesta como suma de exponenciales amortiguadas (Ec.~\ref{eq:response})
    \vspace{0.3em}
    \item \textbf{Discretización:} Mapeo $\lambda_k \to z_k = e^{\lambda_k T_s}$ (Ec.~\ref{eq:mapping})
          $\;\to\;$ recurrencia lineal con polinomio de predicción (Ec.~\ref{eq:ar})
    \vspace{0.3em}
    \item \textbf{Multi-TK:} Hankel multicanal $\to$ SVD $\to$ coeficientes b (Ec.~\ref{eq:Ai}--\ref{eq:bcoeffs})
          $\;\to\;$ polos $\{z_k^{(L)}\}$ como inicialización
    \vspace{0.3em}

    \item \textbf{Patrones dinámicos:} Vandermonde $\to$ residuos $a_{i,k}$ (Ec.~\ref{eq:vander_art}--\ref{eq:residues_art})
          $\;\to\;$ formas modales (Ec.~\ref{eq:mode_shape_art}), factores de participación (Ec.~\ref{eq:pf_mediciones}) y,
          grupos coherentes (Ec.~\ref{eq:direction_cosine_art})
\end{enumerate}

\end{frame}


\begin{frame}
\frametitle{Contenido de los Capítulos 3 }
\framesubtitle{Casos de estudio sintéticos, validación experimental y comparación con otros métodos}


% Columna izquierda: Capítulo 3
\centering
\structure{\textbf{Capítulo 3}} \\
\structure{\small \textit{Estudios sintéticos y pruebas de eficiencia}}

\begin{itemize}
    \item \textbf{Casos sintéticos}
    \begin{itemize}
        \small
        \item Señales multi-componente con parámetros controlados.
        \item Validación de frecuencias, amortiguamientos y amplitudes.
    \end{itemize}
    \item \textbf{Casos base (sistemas de prueba)}
    \begin{itemize}
        \small
        \item Kundur dos áreas.
        \item IEEE 68 barras – 5 áreas.
    \end{itemize}
    \item \textbf{Pruebas de eficiencia del algoritmo}
    \begin{itemize}
        \small
        \item Orden óptimo \(L\) y sensibilidad al umbral de energía.
        \item Costo computacional (tiempo de ejecución, uso de memoria).
        \item Susceptibilidad al ruido (SNR variable).
        \item Comparación de precisión según número de canales.
    \end{itemize}
\end{itemize}

\end{frame}


\begin{frame}
\frametitle{Contenido de los Capítulos 4}
\framesubtitle{Casos de estudio sintéticos, validación experimental y comparación con otros métodos}



\centering
\structure{\textbf{Capítulo 4}} \\
\structure{\small \textit{Caso robusto: LA2030 y señales reales}}

\begin{itemize}
    \item \textbf{Modelo LA2030}
    \begin{itemize}
        \small
        \item Sistema Interconectado Latinoamericano.
        \item Topología, áreas y puntos de medición.
    \end{itemize}
    \item \textbf{Señales reales}
    \begin{itemize}
        \small
        \item Mediciones PMU de eventos reales.
        \item Preprocesamiento: filtrado, diezmado, remuestreo.
    \end{itemize}
    \item \textbf{Comparación con otros métodos}
    \begin{itemize}
        \small
        \item Prony multiseñal.
        \item Matrix Pencil (MP).
        \item ERA (Eigensystem Realization Algorithm).
        \item Métricas: error en frecuencia, amortiguamiento, forma modal.
    \end{itemize}
    \item \textbf{Discusión de resultados}
    \begin{itemize}
        \small
        \item Robustez, precisión y escalabilidad del Multi‑TK.
    \end{itemize}
\end{itemize}

\end{frame}


\ThankYouFrame



\end{document}